% =============================================================================
%  book/ccs.tex  ——  《C++17 By Example》中文版 导言区（Packt 原版版式）
%  参考模板 D:/Ego/trans/latex_source_v2，按 Packt 原书实测版式改写：
%    · 单色（黑/灰），无彩色色带、无盒式代码、无 Note/Tip 盒
%    · 章首页：右侧大号衬线章号 + 无衬线章名（右对齐），无横线
%    · 节标题：粗体无衬线、不编号
%    · 代码：无框、左缩进的等宽字体，强调行加粗；无标题条、无行号
%    · 行内代码：灰色等宽
%    · 目录：扁平缩进层级、无引导点、页码右对齐
%    · 页眉：左斜体章名 + 右 “Chapter N” + 细线；页脚：[ N ] 居中
%    · 索引：两栏无衬线、粗体字母头、等宽页码
% =============================================================================

\documentclass[12bp,openany,oneside,UTF8]{book}
\interfootnotelinepenalty=10000

% -----------------------------------------------------------------------------
%  基础宏包
% -----------------------------------------------------------------------------
\usepackage{textcomp}
\usepackage[T1]{fontenc}
\usepackage{microtype}
\usepackage{amsmath}
\usepackage{amssymb}
\usepackage{ifthen}
\usepackage{pgffor}
\usepackage{xspace}
\usepackage{graphicx}
\usepackage{tikz}
\usepackage{tabularx}
\usepackage{booktabs}
\usepackage{multirow}
\usepackage{float}
\usepackage{enumitem}
\usepackage{indentfirst}
\usepackage{subfiles}
\usepackage{caption}
\usepackage{xcolor}
\usepackage[export]{adjustbox}
\usepackage{multicol}

% -----------------------------------------------------------------------------
%  页面尺寸：Packt 原书实测 540 × 666 bp（7.5in × 9.25in），左右边距约 57/59bp
% -----------------------------------------------------------------------------
\usepackage[
  paperwidth=541.4bp, paperheight=666.1bp,
  left=56.7bp, right=56.7bp,
  top=77.2bp, bottom=68bp,
  headheight=22bp, headsep=16.5bp,
  footskip=19.8bp]{geometry}

% -----------------------------------------------------------------------------
%  字体：中文正文思源宋体；标题/代码注释用微软雅黑（无衬线）；
%        拉丁正文 Cambria（衬线）；代码 Consolas（等宽）
% -----------------------------------------------------------------------------
\usepackage{fontspec}
\usepackage[fontset=none]{ctex}
\usepackage[fontsize=10.5bp]{fontsize}
\linespread{1.0}                         % 正文行距 12.6bp（=1.2×10.5，与原书一致）1.0833(13bp)/1.1250(13.5bp)/1.1667(14bp)/1.2083(14.5bp)/1.2500(15bp)/1.2917(15.5bp)/1.3334(16.0bp)

\setCJKmainfont[AutoFakeSlant=true]{Source Han Serif SC}                 % 中文正文（衬线）
\setCJKsansfont[AutoFakeBold=true,AutoFakeSlant=true]{Microsoft YaHei}   % 中文标题（无衬线）
\setCJKmonofont[AutoFakeBold=true,AutoFakeSlant=true]{Microsoft YaHei}   % 代码块中的中文
\setmainfont{PalatinoLinotype}[
  Path=C:/Windows/Fonts/, Extension=.ttf,
  UprightFont=pala, BoldFont=palab, ItalicFont=palai, BoldItalicFont=palabi]
\setsansfont{Arial}[
  Path=C:/Windows/Fonts/, Extension=.ttf,
  UprightFont=arial, BoldFont=arialbd, ItalicFont=ariali, BoldItalicFont=arialbi]
\setmonofont{Consolas}[
  Path=C:/Windows/Fonts/, Extension=.ttf,
  UprightFont=consola, BoldFont=consolab,
  ItalicFont=consolai, BoldItalicFont=consolaz]
\newfontfamily\codehan{Microsoft YaHei}

% -----------------------------------------------------------------------------
%  颜色（Packt 单色：正文近黑，代码/次要灰）
% -----------------------------------------------------------------------------
\definecolor{ink}{HTML}{1A1A1A}          % 正文墨色（近黑，非纯黑）
\definecolor{headcol}{HTML}{000000}      % 标题（原书黑色）
\definecolor{codegray}{HTML}{444444}     % 行内代码灰
\definecolor{rulecol}{HTML}{000000}      % 页眉细线

\setlength{\parindent}{0bp}              % Packt 正文不缩进，段间留白
\setlength{\parskip}{5.6bp}              % 原书段落间距实测 ~5.7bp
\raggedbottom

\emergencystretch=1em

\AtBeginDocument{\color{ink}}

\setlist[itemize,enumerate,description]{nosep, topsep=4bp, partopsep=0bp, parsep=0bp, leftmargin=2em}

% -----------------------------------------------------------------------------
%  代码高亮：minted（需 Pygments + xelatex -shell-escape）
%  原书代码块无外框/底色/行号，左缩进，强调行加粗
% -----------------------------------------------------------------------------
\usepackage[cache=true]{minted}
\newcommand{\codefont}{\fontsize{9bp}{11.4bp}\selectfont}
\setminted{obeytabs, tabsize=2, breaklines=true, breakanywhere=true,
           fontsize=\codefont, autogobble=true}

% 【关键】minted v3 在 Windows 下缓存目录必须是绝对路径（照搬模板做法）
\makeatletter
\CatchFileDef{\mintedProjectRoot}{|"kpsewhich --var-value TEXMF_OUTPUT_DIRECTORY"}%
  {\minted@standardcatcodes\endlinechar=-1}%
\def\minted@stripcr#1\r{#1}%
\edef\mintedProjectRoot{\expandafter\minted@stripcr\mintedProjectRoot\r}%
\ifx\mintedProjectRoot\@empty
  \CatchFileDef{\mintedProjectRoot}{|"kpsewhich --var-value TEXMFOUTPUT"}%
    {\minted@standardcatcodes\endlinechar=-1}%
  \edef\mintedProjectRoot{\expandafter\minted@stripcr\mintedProjectRoot\r}%
\fi
\ifx\mintedProjectRoot\@empty
  \PackageWarning{minted-cn}{Cannot determine project root...}%
\else
  \edef\minted@cachedir{\mintedProjectRoot}%
  \typeout{minted-cn: cachedir = \minted@cachedir}%
\fi
\makeatother

\makeatletter
\CatchFileDef{\minted@styledef@manning}{book/manning-pygments-style.tex}%
  {\minted@standardcatcodes\endlinechar=-1}%
\minted@detectconfig
\makeatother
\setminted{style=manning}
\usepackage[most]{tcolorbox}
\tcbuselibrary{breakable, minted, skins}

% 无框代码块：左缩进由 \codeInset 控制（原书多数章节 63bp；第 3–5 章 18bp）
\newcommand{\codeInset}{63bp}
\newtcblisting{codebox}{
  listing engine=minted, minted language=c++, listing only,
  size=minimal, arc=0bp, breakable, enhanced jigsaw,
  colframe=white, colback=white, boxrule=0bp,
  left=\codeInset, right=0bp, top=0bp, bottom=0bp, boxsep=0bp,
  before skip=9.4bp, after skip=9.7bp,
  minted options={fontsize=\codefont, tabsize=2, obeytabs=true,
                  breaklines=true, breakanywhere=true, autogobble=true},
}
% 带标题代码块（原书其实无标题条；此环境保留以备个别需要标注文件名处）
\newtcblisting{codeboxt}[1]{
  listing engine=minted, minted language=c++, listing only,
  size=title, arc=0bp, breakable, enhanced jigsaw,
  colframe=white, colback=white, boxrule=0bp,
  left=15bp, right=0bp, top=2bp, bottom=2bp, boxsep=0bp,
  lefttitle=12bp, righttitle=2bp, toptitle=2bp, bottomtitle=2bp,
  title={#1},
  colbacktitle=white, coltitle=ink,
  fonttitle=\sffamily\fontsize{9bp}{11bp}\selectfont,
  before skip=8bp, after skip=8bp,
  minted options={fontsize=\codefont, tabsize=2, obeytabs=true,
                  breaklines=true, breakanywhere=true, autogobble=true},
}

% -----------------------------------------------------------------------------
%  Note / Tip 提示块（原书实测版式，pdf 17/28/65/171/179/318）
%    图标 41.5×41.5bp 灰色圆（i / TIP），图片左缘 x=71.3（正文左缘 +14.6bp）；
%    正文整块缩进 65.4bp（自 x=56.7 到 x=122.1），无加粗标题行、无底色盒子。
%  图标纵向位置（实测反推的原始版式规则）：
%    图标上缘 = 块顶 + max(3.6bp, (块高 − 41.5bp)/2)
%    等价为：图标上缘相对首行基线 delta = max(−5.8bp, (块高 −41.5bp)/2 −10bp)
%    短块（1–2 行）图标上缘在首行基线上方 5.8bp；6 行块图标在块内基本居中（误差 <0.7bp）。
% -----------------------------------------------------------------------------
% \myNotic{图标文件名}{正文}，图标文件名如 p017_i00（i）/ p017_i01（TIP）
% 纵向规则（原书实测反推）：
%   多行块（>2 行）：图标与正文块的垂直中心对齐；
%   短块（≤2 行）：图标上缘固定在首行基线上方 5.8bp（与原书 pdf 17 一致）。
% 用 \smash 把图标包成“零高零深”，避免撑高行盒、破坏段间距。
\newdimen\noticHT
\newdimen\noticR\newdimen\noticPad
\newbox\noticBox
\newenvironment{myNotic}[2]{%
  \par\addvspace{4.6bp}%
  \setbox\noticBox\vbox{\hsize=\dimexpr\textwidth-65.4bp\relax
    \parindent 0bp \parskip 0bp \raggedright \noindent #2\par}%
  \noticHT=\ht\noticBox \advance\noticHT by \dp\noticBox
  % 图标纵向定位（实测标定：图标上缘 ≈ 首行上缘 +2.0bp，如原书 pdf 17/65/171/179）
  %   短块（≤2 行）：R = −34.2bp
  %   多行块：图标与正文块垂直居中，中心略高 1.6bp（原书 pdf 28/318）→ R = −块高/2 − 10.6bp
  \ifdim \noticHT > 28bp
    \noticR=-\noticHT
    \divide\noticR by 2
    \advance\noticR by -10.6bp
  \else
    \noticR=-34.2bp
  \fi
  \noindent\hangindent 65.4bp\hangafter 1\relax
  \rlap{\hskip 14.6bp\smash{\raisebox{\noticR}{\includegraphics[width=41.5bp,height=41.5bp]{images/#1.png}}}}%
  \hspace*{65.4bp}%
  \parindent 0bp\parskip 0bp\relax
  #2\par
  % 图标高于正文块时块后补足（原书 pdf 17 实测：两相邻提示块正文上缘间距 59.4bp）
  \noticPad=41.5bp \advance\noticPad by -\noticHT \advance\noticPad by 5.3bp
  \ifdim\noticPad<0bp \noticPad=0bp \fi
  \vspace{\noticPad}\par\addvspace{4.6bp}}

% 补丁：minted v3 × tcolorbox × subfiles 路径冲突（照搬模板）
\makeatletter
\def\minted@define@inputfilepath#1{%
  \IfPackageLoadedTF{import}%
   {\IfFileExists{\import@path#1}%
     {\edef\minted@inputfilepath{\import@path#1}}%
     {\edef\minted@inputfilepath{#1}}}%
   {\edef\minted@inputfilepath{#1}}}
\makeatother

% -----------------------------------------------------------------------------
%  图表
% -----------------------------------------------------------------------------
\newif\ifmyGraphicCapCenter
\myGraphicCapCenterfalse
\newif\ifmyGraphicCapBold
\myGraphicCapBoldfalse            % Packt 图注不加粗
\ExplSyntaxOn
\NewDocumentCommand { \myGraphic } { mmm }
{
  \par\medskip\noindent
  \begin{center}
    \includegraphics[width=#1\textwidth,height=0.82\textheight,keepaspectratio]{#2}
  \end{center}
  \nopagebreak\vspace{3bp}\noindent
  {\fontsize{8bp}{10.5bp}\selectfont
   \ifmyGraphicCapBold\bfseries\fi
   \ifmyGraphicCapCenter
     \parbox{\textwidth}{\centering #3}%
   \else
     \parbox{\textwidth}{#3}%
   \fi}\par
  \medskip
}
\ExplSyntaxOff
\captionsetup[figure]{labelformat=empty}
\captionsetup[table]{labelformat=empty}

% -----------------------------------------------------------------------------
%  标题样式（Packt：粗体无衬线、不编号；黑色）
% -----------------------------------------------------------------------------
\usepackage{titlesec}
\titleformat{\section}
  {\sffamily\bfseries\fontsize{20bp}{24bp}\selectfont\color{headcol}}{}{0pt}{}
\titlespacing*{\section}{0pt}{20.9bp}{-4.5bp}
\titleformat{\subsection}
  {\sffamily\bfseries\fontsize{18bp}{22bp}\selectfont\color{headcol}}{}{0pt}{}
\titlespacing*{\subsection}{0pt}{16.6bp}{-4.4bp}
\titleformat{\subsubsection}
  {\sffamily\bfseries\fontsize{16bp}{20bp}\selectfont\color{headcol}}{}{0pt}{}
\titlespacing*{\subsubsection}{0pt}{10bp}{4bp}

\setcounter{secnumdepth}{0}     % 不显示节号（对齐原书）
\setcounter{tocdepth}{2}        % 目录含章/节/小节（原书目录为三级）

% -----------------------------------------------------------------------------
%  目录 / 书签 / 超链接
% -----------------------------------------------------------------------------
\usepackage{url}
\usepackage[depth=3]{bookmark}
\usepackage{hyperref}
\hypersetup{
  colorlinks=true,
  linkcolor=black,
  citecolor=black,
  filecolor=black!35!blue,
  urlcolor=black!35!blue,
  breaklinks=true,
  pdfstartview=Fit,
  pdftitle={C++17 边做边学},
  pdfauthor={Stefan Bjornander},
  pdfsubject={C++17 By Example 中文版},
  pdfkeywords={C++, C++17, Qt, DSL},
}
\renewcommand{\contentsname}{目录}

% 目录条目：三级（章粗体 / 节粗体缩进 / 小节常规缩进），无引导点，页码右对齐
\makeatletter
% 原书目录实测：章条 x=59.6 粗体 12；节条 x=77.6 粗体 12；小节条 x=95.6 常规 11；
% 页码右缘 x=481.8；条距 章 4.2bp / 节 1.2bp / 小节 1.0bp（正文 parskip 置零）
\renewcommand*\l@chapter[2]{%
  \ifnum \c@tocdepth >\m@ne
    \addpenalty{-\@highpenalty}%
    \parskip\z@
    \vskip 4.2bp \@plus\p@
    \begingroup
      \leavevmode \bfseries\sffamily\fontsize{12bp}{14bp}\selectfont
      \leftskip 2.9bp\relax
      #1\nobreak\hfill\nobreak\hb@xt@\@pnumwidth{\hss #2\hspace*{2.9bp}}\par
      \penalty\@highpenalty
    \endgroup
  \fi}
\renewcommand*\l@section[2]{%
  \ifnum \c@tocdepth >\z@
    \parskip\z@
    \vskip 1.2bp \@plus\p@
    \begingroup
      \leavevmode \bfseries\sffamily\fontsize{12bp}{14bp}\selectfont
      \leftskip 20.9bp\relax
      #1\nobreak\hfill\nobreak\hb@xt@\@pnumwidth{\hss #2\hspace*{2.9bp}}\par
    \endgroup
  \fi}
\renewcommand*\l@subsection[2]{%
  \ifnum \c@tocdepth >\z@
    \parskip\z@
    \vskip 1.0bp \@plus\p@
    \begingroup
      \leavevmode \sffamily\fontsize{11bp}{13bp}\selectfont
      \leftskip 38.9bp\relax
      #1\nobreak\hfill\nobreak\hb@xt@\@pnumwidth{\hss #2\hspace*{2.9bp}}\par
    \endgroup
  \fi}
% 前置/后置页条目下方的整宽细线（原书 Preface / 书末 / 索引条下各一条）
\newcommand{\tocentryrule}{%
  \par\vspace*{-10.3bp}\noindent\rule{\textwidth}{0.4pt}\par\vspace{-1.7bp}}

% 自定义 \tableofcontents：右对齐“目录”标题（基线 104.9）+ 首页无页眉
% 原书首页：标题 baseline 104.9，首条 Preface baseline 149.5（其下 151.9 处整宽细线）
\renewcommand{\tableofcontents}{%
  \clearpage
  \pagenumbering{roman}%
  \thispagestyle{empty}%
  \vspace*{-23.3bp}%
  \begin{flushright}
    {\sffamily\fontsize{30bp}{34bp}\selectfont\contentsname}
  \end{flushright}
  \vspace{15.5bp}%
  \markboth{}{\contentsname}%  目录续页页眉：右侧“目录”（与原书一致）
  \@starttoc{toc}%
}
\makeatother

% -----------------------------------------------------------------------------
%  长串断行（URL / 行内长标识符允许折行，不补连字符）
% -----------------------------------------------------------------------------
\ExplSyntaxOn
\let\myORIhref\href
\let\myORItt\texttt
\cs_new:Npn \my_spcmark {}
\tl_new:N \l_myin_tl
\tl_new:N \l_myout_tl
\cs_new_protected:Npn \my_breakinject:n #1
  {
    \tl_set:Nn \l_myin_tl { #1 }
    \tl_replace_all:Nnn \l_myin_tl { ~ } { \my_spcmark }
    \tl_clear:N \l_myout_tl
    \tl_map_inline:Nn \l_myin_tl
      { \tl_put_right:Nn \l_myout_tl { ##1 \penalty 1000 \scan_stop: } }
    \tl_replace_all:Nnn \l_myout_tl { / \penalty 1000 \scan_stop: } { / \penalty 0 \scan_stop: }
    \tl_replace_all:Nnn \l_myout_tl { . \penalty 1000 \scan_stop: } { . \penalty 0 \scan_stop: }
    \tl_replace_all:Nnn \l_myout_tl { - \penalty 1000 \scan_stop: } { - \penalty 0 \scan_stop: }
    \tl_replace_all:Nnn \l_myout_tl { ? \penalty 1000 \scan_stop: } { ? \penalty 0 \scan_stop: }
    \tl_replace_all:Nnn \l_myout_tl { = \penalty 1000 \scan_stop: } { = \penalty 0 \scan_stop: }
    \tl_replace_all:Nnn \l_myout_tl { \_ \penalty 1000 \scan_stop: } { \_ \penalty 0 \scan_stop: }
    \tl_replace_all:Nnn \l_myout_tl { \my_spcmark } { ~ }
  }
% 行内代码着灰色
\RenewDocumentCommand { \href } { m m }
  { \my_breakinject:n { #2 } \myORIhref { #1 } { \l_myout_tl } }
\RenewDocumentCommand { \texttt } { m }
  { \my_breakinject:n { #1 } \textcolor{codegray}{\myORItt { \l_myout_tl }} }
\ExplSyntaxOff

% -----------------------------------------------------------------------------
%  页眉页脚（Packt：左斜体章名 + 右 “Chapter N” + 细线；
%            页脚 页码两侧各有细线，页码 12bp 加粗居中）
% -----------------------------------------------------------------------------
\makeatletter
\renewcommand*\@pnumwidth{2.2em}
\renewcommand*\@tocrmarg{2.2em}
\makeatother

\usepackage{fancyhdr}
\pagestyle{fancy}
\fancyhf{}
\renewcommand{\headrulewidth}{0.4pt}
\renewcommand{\headrule}{\hbox to\headwidth{\color{rulecol}\leaders\hrule height 0.4pt\hfill}}
\renewcommand{\footrulewidth}{0pt}
% 页脚：页码居中、两侧各一段细线（与原书一致，线穿过页码中线）
\newcommand{\myfoot}[1]{%
  \leavevmode\hbox to\headwidth{%
    \color{rulecol}%
    \rule[3.6bp]{192.6bp}{0.4pt}%
    \hskip 11.4bp
    \hbox to 20bp{\hss\color{ink}\fontsize{12bp}{14bp}\selectfont\bfseries #1\hss}%
    \hskip 11.4bp
    \rule[3.6bp]{192.6bp}{0.4pt}%
  }}
\fancyhead[L]{\raisebox{-7.1bp}[0bp][0bp]{\itshape\color{ink}\fontsize{10bp}{12bp}\selectfont\leftmark}}
\fancyhead[R]{\raisebox{-7.1bp}[0bp][0bp]{\itshape\color{ink}\fontsize{10bp}{12bp}\selectfont\rightmark}}
\fancyfoot[C]{\myfoot{[~\thepage~]}}
\fancypagestyle{plain}{%
  \fancyhf{}%
  \fancyhead[L]{\raisebox{-7.1bp}[0bp][0bp]{\itshape\color{ink}\fontsize{10bp}{12bp}\selectfont\leftmark}}%
  \fancyhead[R]{\raisebox{-7.1bp}[0bp][0bp]{\itshape\color{ink}\fontsize{10bp}{12bp}\selectfont\rightmark}}%
  \renewcommand{\headrulewidth}{0.4pt}%
  \fancyfoot[C]{\myfoot{[~\thepage~]}}}

% 章节命令（Packt 版式：右侧 100bp 章号 + 30bp 章名，无页眉页脚）
\ExplSyntaxOn
% 编号正文章节：\myChapter{章号}{标题}{文件}
\NewDocumentCommand { \myChapter } { mmm }
{
  \clearpage
  \setcounter{chapter}{#1}
  \setcounter{section}{0}
  \setcounter{subsection}{0}
  \thispagestyle{empty}
  \vspace*{-6.9bp}
  \begin{flushright}
    {\rmfamily\fontsize{100bp}{100bp}\selectfont #1}\par
    \vspace{4bp}
    {\sffamily\fontsize{30bp}{34bp}\selectfont #2}
  \end{flushright}
  \vspace{42bp}
  \markboth{#2}{Chapter~#1}
  \addcontentsline{toc}{chapter}{Chapter~#1:\ \ #2}
  \subfile{#3}
}
% 前言类：进目录但无编号（右对齐 30bp 标题，首页无页眉页脚）
% 标题基线 = 104.9bp（原书实测）：用 \vspace* 拉到 (104.9 - 32bp) 的顶端位置
\NewDocumentCommand { \myFront } { mm }
{
  \clearpage
  \thispagestyle{empty}
  \vspace*{-24.1bp}
  \begin{flushright}
    {\sffamily\fontsize{30bp}{34bp}\selectfont #1}
  \end{flushright}
  \vspace{30bp}
  \markboth{#1}{}
  \addcontentsline{toc}{chapter}{#1}
  \addtocontents{toc}{\protect\tocentryrule}%  前置/后置页条目下方的整宽细线
  \subfile{#2}
}
% 扉页/版权类：不进目录
\NewDocumentCommand { \myFrontNoToc } { m }
{
  \clearpage
  \markboth{}{}
  \subfile{#1}
}
\ExplSyntaxOff

% 节标题不改动页眉标记（保持右栏恒为 “Chapter N”）
\renewcommand{\sectionmark}[1]{}

% -----------------------------------------------------------------------------
%  Unicode 符号兜底
% -----------------------------------------------------------------------------
\usepackage{newunicodechar}
\newunicodechar{→}{\(\to\)}
\newunicodechar{✓}{\(\checkmark\)}
\newunicodechar{ϵ}{\(\epsilon\)}
\newunicodechar{│}{\(|\)}
\newunicodechar{▶}{\(\triangleright\)}
\newunicodechar{·}{\(\cdot\)}

% 索引支持（两栏，粗体字母头 16.2bp，词条 Arial 9bp，页码 FreeMono 9bp）
% 原书实测：字母头 baseline→首条 13.5bp；条目行距 11.5bp；子条目缩进 7.5bp
\newfontfamily\idxnumfont[
  Path=D:/ProgramData/texlive/2026/texmf-dist/fonts/opentype/public/gnu-freefont/,
  Extension=.otf, UprightFont=FreeMono,
  BoldFont=FreeMonoBold, ItalicFont=FreeMonoOblique]{FreeMono}
% 字母头（原书实测：上一组末条 → 字母头 23.4bp；字母头 → 本组首条 13.5bp）
\newcommand{\idxletter}[1]{%
  \par\vspace{9.6bp}\noindent
  {\sffamily\bfseries\fontsize{16.2bp}{16.2bp}\selectfont #1}\par\vspace{2.0bp}}% 折行续行缩进（原书实测：条目首行 x=56.7，续行 x=70.2 即悬挂缩进 13.5bp）
\newcommand{\idxitem}[2]{%
  {\par\noindent\hangindent 13.5bp\hangafter 1\relax
   \sffamily\fontsize{9bp}{11.5bp}\selectfont #1\hskip 4bp
   {\idxnumfont #2}\par}}
\newcommand{\idxsub}[2]{%
  {\par\noindent\hangindent 13.5bp\hangafter 1\relax\hskip 7.5bp
   \sffamily\fontsize{9bp}{11.5bp}\selectfont #1\hskip 4bp
   {\idxnumfont #2}\par}}
\newenvironment{myindex}
  {\par\parindent 0bp\relax\parskip 0bp\relax\frenchspacing
   \begin{multicols}{2}\setlength{\columnsep}{14.2bp}
   \parindent 0bp\relax\parskip 0bp\relax}
  {\end{multicols}}
